\FloatBarrier
\subsection{Whole-Scene Multi-Method Comparisons}
\label{sec:kitti-scene-visual}

\thesisfigure{k_scene_A}{Line A multi-method scene comparison}{
Lab frame 000590, historical Line~A E1 inputs. Gray points are retained
observed rows and blue points are generated rows. The shared crop is
0--60~m forward and $-18$--18~m lateral, all heights, with no display
subsampling. PU-Net* is the old faulty adapter; PU-EF is the lab
compatibility path.}{fig:k-scene-A}

\thesisfigure{k_scene_B}{Line B multi-method scene comparison}{
The same scene under historical Line~B construction. Measured points
now come from the sparse baseline. All four stored method outputs are
included under the same crop and colour convention as
Figure~\ref{fig:k-scene-A}.}{fig:k-scene-B}

Figures~\ref{fig:k-scene-A} and~\ref{fig:k-scene-B} re-render the local
lab visualisation arrays for frame 000590, including PDANS, PU-GCN,
PU-EdgeFormer, and PU-Net in both lines. The display window is fixed to
0--60~m forward and $-18$--18~m lateral, with no height filter.
All points inside this window are drawn. Measured rows are identified
by exact float32-row membership in the corresponding observed input,
rather than by assigning a method a presumed preservation property.

The separate colours reveal a distinction hidden by an all-blue dense
cloud: these historical E1 inputs retain observations at the stored-cloud
stage, but much of the displayed density comes from generated rows.
In Line~B, the measured component is the sparse input, so recovery
to the original file size does not restore the original measurements.
The figure's point counts refer to the display window, not to detector
input cardinality or the number of points inside labelled objects.

\FloatBarrier
\subsection{Local Surfaces and Measurement Preservation}
\label{sec:kitti-local-visual}

\thesisfigure{k_local_A}{Line A local multi-method surfaces}{
Common object-aligned crop for frame 005625, Car audit index 10.
All four outputs retain 316 measured rows in this window; generated
rows are shown separately. The dashed GT box and view are identical
across methods. Historical pipeline labels match Figure~\ref{fig:k-scene-A}.}{fig:k-local-A}

\thesisfigure{k_local_B}{Line B local multi-method surfaces}{
The same labelled Car after sparsification. All points within the
declared crop are shown; counts distinguish the 85 measured rows from
generated rows. Increased crop density alone does not establish
restored surface accuracy or improved detection.}{fig:k-local-B}

The whole-scene view cannot resolve the distribution around one object.
Figures~\ref{fig:k-local-A} and~\ref{fig:k-local-B} therefore use the
same Car in frame 005625 (audit GT index 10) and the same object-aligned
coordinates for every method. The crop extends 0.6~m beyond the box in
length and width and 0.4~m in height. These margins retain context
around the measured surface; crop counts are not restricted to the
interior of the ground-truth box.

The original crop contained 316 observed rows, compared with 85 after
sparsification. In Line~A, PDANS added 1,184 generated rows within this
window and PU-GCN added 953. In Line~B, the respective counts were
284 and 259. The visible point concentration and the incomplete
measurement support can therefore be assessed separately.
The dashed box specifies the labelled object's extent; it is not a
surface reconstruction target. Points outside it can be legitimate
background because the crop deliberately includes a margin.

\FloatBarrier
\subsection{Matched Detection Failures and Positive Counterexamples}
\label{sec:kitti-matched-cases}

\thesisfigure{k_case_pr_lost}{PointRCNN Car loss case}{Historical lab frame 000590, Line~A, PDANS, audit object index 2. The oblique view and BEV show the same effective-input crop. Predictions are read from the retained manifest; only qualifying matches are drawn. The class threshold is 0.70 3D IoU.}{fig:k-case-pr-lost}

The four examples below were already identified in the lab event audit.
They are selected diagnostic cases, not a random estimate of event
frequencies. Each compares the same labelled object before and after
upsampling using detector-effective input. Dashed boxes are ground
truth; solid boxes are saved qualifying predictions. Unqualified
predictions are not drawn, so an FN label means no qualifying match,
not necessarily an absence of proposals.

The PointRCNN Line~A PDANS example lost a Car that initially had a well-aligned prediction. The common local view permits inspection of altered neighbourhood support without confusing this with the global increase in stored point count.
Figure~\ref{fig:k-case-pr-lost} records 3D IoU 0.777 before upsampling and no qualifying match afterwards. The displayed effective-input crop contains 846 baseline points and 785 upsampled points.

\FloatBarrier

A different PointRCNN example provides a positive counterexample in Line~B. The recovered Car crosses the matching threshold despite fewer effective input points in this local crop. Local cardinality alone therefore cannot explain the sign of a detection transition.
Figure~\ref{fig:k-case-pr-recovered} records no qualifying match before upsampling and 3D IoU 0.705 afterwards. The displayed effective-input crop contains 312 baseline points and 271 upsampled points.

\thesisfigure{k_case_pr_recovered}{PointRCNN Car recovery case}{Historical lab frame 005625, Line~B, PDANS, audit object index 7. The oblique view and BEV show the same effective-input crop. Predictions are read from the retained manifest; only qualifying matches are drawn. The class threshold is 0.70 3D IoU.}{fig:k-case-pr-recovered}
\FloatBarrier

The CenterPoint Line~B PU-GCN Car is a complementary negative case: more retained crop points coexist with loss of a previously qualifying prediction. This bounds the claim that preserving measurements or adding local density is sufficient for detection.
Figure~\ref{fig:k-case-cp-lost} records 3D IoU 0.871 before upsampling and no qualifying match afterwards. The displayed effective-input crop contains 100 baseline points and 394 upsampled points.

\thesisfigure{k_case_cp_lost}{CenterPoint Car loss case}{Historical lab frame 005625, Line~B, PU-GCN, audit object index 10. The oblique view and BEV show the same effective-input crop. Predictions are read from the retained manifest; only qualifying matches are drawn. The class threshold is 0.70 3D IoU.}{fig:k-case-cp-lost}
\FloatBarrier

The CenterPoint Cyclist in the same Line~B PU-GCN frame provides a positive case under the class-specific 0.50 IoU threshold. It shows that recoveries can occur in the same representation and scene as Car losses; the full-validation event counts establish which direction dominates.
Figure~\ref{fig:k-case-cp-recovered} records no qualifying match before upsampling and 3D IoU 0.655 afterwards. The displayed effective-input crop contains 19 baseline points and 110 upsampled points.

\thesisfigure{k_case_cp_recovered}{CenterPoint Cyclist recovery case}{Historical lab frame 005625, Line~B, PU-GCN, audit object index 4. The oblique view and BEV show the same effective-input crop. Predictions are read from the retained manifest; only qualifying matches are drawn. The class threshold is 0.50 3D IoU.}{fig:k-case-cp-recovered}
\FloatBarrier

These cases discriminate several explanations that an AP table cannot.
A global fourfold expansion need not increase the number of points
entering a particular PointRCNN crop, and a larger CenterPoint crop
can still lose its matched box. Conversely, a positive transition does
not require more local effective points. The examples support analysing
where generated points enter the representation and how predictions
change, while leaving the causal contribution of geometry, sampling,
and learned features unresolved. The full-validation final matrix
remains the evidence for detector adaptation.
